Classical variance reduction transfers to diffusion training in the sense that the estimators are unbiased,
cheap and act as theory predicts. In this regime they do not translate into better models at equal compute.
The law-of-total-variance decomposition explains why. The DDPM gradient's variance is dominated by data and
noise \emph{within} a timestep, not by the choice of timestep, which caps what stratification can achieve. Antithetic
noise cancels the noise term efficiently, but pays for it by halving the number of distinct images. The finance
lesson that transfers is diagnostic: measure the variance components first, then choose the estimator. The
most promising direction the measurements point to is a gradient-norm-based importance proposal, measured
post hoc at $\approx$\PHOracle$\times$ the baseline variance at fixed parameters, combined with stratification.
The only training-outcome signal, a 1.5--3\% held-out-loss gain from loss-based importance sampling, is
consistent across both datasets but not conclusive. That is a conjecture to test at larger
scale and with more seeds.
